\section{The primary-vertex bias and the modified reconstruction}
\label{sec:fixes}

\subsection{How the stock reconstruction uses the PV}
The \Dzero{} instance of \texttt{KFParticle\_sPHENIX} (\texttt{[D0 -> K\^{}- pi\^{}+]cc}) is configured with
\fcmd{dontUseGlobalVertex(true)} and \fcmd{constrainToPrimaryVertex(true)}: it reads the \texttt{SvtxVertexMap}, and for
every track pair and mass assignment it tries every PV of the event, keeps the combinations that pass the mother cuts,
and among those keeps one per track pair, the one with the smallest PV--DCA $\chi^2$ (\fcmd{selectBestCombination}, read
from source). Pointing (DIRA), decay length and its significance are computed w.r.t.\ that PV. The track--PV DCA cut of
the daughters uses the minimum DCA over the PVs of the track's crossing (\fcmd{calcMinPV_DCA}).

The \texttt{SvtxVertexMap} vertices are found by \texttt{PHSimpleVertexFinder} from the tracks of each crossing; the
\Dzero{} daughters are among them. After vertex finding, \texttt{PHActsVertexPropagator} moves every track state to its
point of closest approach to its vertex. This was checked in a DST (\fpath{tools/pvfix/trkstate_check.C},
\tsVtxTracks{} vertex tracks): the component of (state $-$ vertex) along the track momentum has a median of
\tsAlong\um. The states the vertex finder used are therefore no longer in the DST.

\subsection{Changes}
\begin{description}
\item[Fix 1 (bookkeeping).] The ntuple branch \texttt{primary\_vertex\_ID} held the beam crossing, not the
  \texttt{SvtxVertexMap} key; it now holds the key, and the crossing is in \texttt{primary\_vertex\_crossing}.
\item[Fix 2 (pair refit).] For each candidate the tracks attached to its PV are selected (quality $\leq 300$,
  $\pt\geq0.1\gevc$, $\geq3$ MVTX clusters on the silicon seed) and the vertex is recomputed with a re-implementation of
  \texttt{PHSimpleVertexFinder}: straight-line points of closest approach for all track pairs with DCA $<500\um$
  (if none, $<1.5$\,mm) and both points within $|x|,|y|<2$\,mm; the median of all points; the average of the pairs within
  1\,mm of the median in $x$ and $y$. Covariance: sum of the used tracks' position covariances divided by
  $\sqrt{N}$. This is done twice: with all selected tracks (\emph{refitAll}, closure test) and without the candidate's
  daughters (\emph{refit}).
\item[Fix 3 (KF fit).] The same track lists are fitted with \fcmd{KFVertex::ConstructPrimaryVertex}, the tracks entering
  with their full DST state and $6\times6$ covariance (as \fcmd{KFParticle_Tools::makeParticle}); \emph{kfAll} and
  \emph{kf}. Fixes 3b/3c make the outlier cut and a covariance scale configurable (Sec.~\ref{sec:splithalf}).
\item[Fix 4 (selection with any PV).] Flags \texttt{pvfix\_passOriginal/Refit/KF} record whether the candidate passes
  the production cuts (both daughters $|{\rm DCA}|\geq30\um$, DIRA $\geq0.9$) w.r.t.\ each PV; \texttt{passOriginal}
  reproduces the stock in-job cut (minimum DCA over the PVs of the crossing). With the in-job cuts loosened, a candidate
  is written if any flag is set.
\item[Fix 5 (stock combination preserved).] With a looser in-job DIRA cut, a combination that fails the stock DIRA cut
  could win the best-combination choice and replace the stock candidate (\smNoFiveMissB{} of \smStockB{} stock candidates were lost in a
  smoke test of run \smRunB). The choice is now restricted to combinations with DIRA $\geq 0.9$ whenever one exists.
\end{description}
For every PV version the ntuple stores the vertex position, status and number of tracks, and the candidate quantities
(DIRA, DIRA$_{xy}$, flight-distance $\chi^2$, PV--DCA, decay length and its error, in 3D and $xy$; per daughter DCA,
DCA$_{xy}$ and their significances) with the suffixes \texttt{\_PVrefit} and \texttt{\_PVkf}.

\subsection{KF-fit track errors from split-half vertices}
\label{sec:splithalf}
The track covariances in the DST turned out to be too small for vertexing. The test (\fpath{tools/pvfix/pv_splithalf.C},
\splitVtx{} vertices from one DST file each of runs 79513 and 80880, \splitSel{} with $\geq\splitMinTrk$ selected tracks) splits the tracks of each
vertex by index parity, fits each half with each method, and uses $\sigma(\Delta)/2$ of the two half-vertices as the
resolution of the full vertex; for the KF fit, $\Delta/\sqrt{C_1+C_2}$ tests the error scale. Widths are robust
(1.4826\,$\times$\,MAD).

\begin{table}[H]
\centering\small
\caption{Split-half test of the PV methods. Scale: factor applied to the track covariance (squared); resolution:
$\sigma(\Delta)/2$; pull: width of $\Delta/\sqrt{C_1+C_2}$; used: median fraction of tracks kept by the outlier cut.
Source: \fpath{logs/split_s1.txt}, \fpath{logs/split_s7.txt}.}
\label{tab:splithalf}
\begin{tabular}{@{}l c r c c c c@{}}
\toprule
method & scale & vertices & \hd{resolution $x$ / $z$}{(\um)} & pull $x$ / $z$ & $\chi^2$/ndf & used \\
\midrule
\csname @@input\endcsname tables/splithalf.tex 
\bottomrule
\end{tabular}
\end{table}

With the DST errors the KF fit has $\chi^2/{\rm ndf}=\spKallOneChi$ and pulls of \spKallOnePullX{} ($x$) and
\spKallOnePullZ{} ($z$); the default outlier cut (3.5) then rejects most tracks (\spKcutOneUse{} kept) and only
\spKcutOneN{} of \splitSel{} vertices survive. Scaling the errors by 7 gives $\chi^2/{\rm ndf}=\spKallSevenChi$ and
pulls \spKallSevenPullX/\spKallSevenPullZ{}; with the cut at 3.5 the resolution is \spKcutSevenX/\spKcutSevenZ\um{}
($x$/$z$), against \spPairOneX/\spPairOneZ\um{} for the pair average.

\choice{KF fit with track errors $\times7$ and outlier cut 3.5. The scale is calibrated on the data themselves: without
outlier rejection it brings the vertex $\chi^2$/ndf from \spKallOneChi{} to \spKallSevenChi{} and the pulls from
\spKallOnePullX/\spKallOnePullZ{} to \spKallSevenPullX/\spKallSevenPullZ; with the cut at 3.5 it gives the best
split-half resolution of the methods tested. The $z$ pull is then below 1 (\spKcutSevenPullZ), i.e.\ the $z$ errors are
somewhat overestimated.}
\caveat{The scaled errors also enter the vertex covariance and hence $\sigma_L$: with the KF fit, decay-length
significances are smaller than with the other PVs, and the same \Ls{} cut selects far fewer candidates
(Table~\ref{tab:counts}). A \Ls{} cut is therefore not equivalent across PV versions. The factor 7 is a single number for
all tracks; its \pt{} and $\eta$ dependence was not studied.}
